% Exported from Math++. Compile with pdflatex.
\documentclass[11pt,a4paper]{article}
\usepackage[T1]{fontenc}
\usepackage[utf8]{inputenc}
\usepackage{lmodern}
\usepackage{amsmath,amssymb,amsthm,mathtools}
\usepackage{graphicx,booktabs,xcolor}
\usepackage[hidelinks]{hyperref}
\usepackage[margin=25mm]{geometry}
\providecommand{\degree}{\ensuremath{^\circ}}
\providecommand{\xmark}{\ensuremath{\times}}
\providecommand{\omicron}{o}
\providecommand{\oiint}{\mathop{\oint\mkern-12mu\oint}}
\theoremstyle{plain}
\newtheorem{theorem}{Theorem}
\newtheorem{lemma}[theorem]{Lemma}
\newtheorem{corollary}[theorem]{Corollary}
\newtheorem{proposition}[theorem]{Proposition}
\newtheorem{conjecture}[theorem]{Conjecture}
\theoremstyle{definition}
\newtheorem{definition}[theorem]{Definition}
\newtheorem{example}[theorem]{Example}
\theoremstyle{remark}
\newtheorem*{remark}{Remark}
\setlength{\parindent}{0pt}\setlength{\parskip}{0.6em}
\begin{document}
\section*{Limits}

\begin{equation*}
  \lim_{x \to 0} \frac{\sin x}{x} = 1
\end{equation*}
\begin{equation*}
  \lim_{n \to \infty} \left(1 + \frac{1}{n}\right)^n = e
\end{equation*}
\begin{equation*}
  \lim_{x \to \infty} \frac{3 x^2 + 1}{x^2 - 5} = 3
\end{equation*}
\begin{equation*}
  \lim_{x \to 0} \frac{1 - \cos x}{x^2} = \frac{1}{2}
\end{equation*}
\begin{equation*}
  \lim_{x \to 2} \frac{x^2 - 4}{x - 2} = 4
\end{equation*}
\end{document}
